\documentclass[9pt,a4paper,twocolumn]{ctexart}

% --- Packages ---
\usepackage[top=1.2cm, bottom=1.2cm, left=1.2cm, right=1.2cm, columnsep=0.8cm]{geometry}
\usepackage{hyperref}
\usepackage{graphicx}
\usepackage{float}
\usepackage{longtable}
\usepackage{tabularx}
\usepackage{booktabs}
\usepackage{enumitem}
\usepackage{xcolor}
\usepackage{fancyhdr}
\usepackage{listings}
\usepackage{fontspec}
\usepackage{titlesec}
\usepackage{caption}

% --- Code block styling ---
\definecolor{codebg}{RGB}{245,245,245}
\definecolor{codeframe}{RGB}{220,220,220}
\definecolor{codekw}{RGB}{0,0,192}
\definecolor{codecomment}{RGB}{0,128,0}
\definecolor{codestring}{RGB}{163,21,21}
\definecolor{codenum}{RGB}{128,0,128}
\definecolor{codepunct}{RGB}{90,90,90}
\definecolor{badgegreen}{RGB}{34,139,34}
\definecolor{badgeblue}{RGB}{30,144,255}
\definecolor{badgeblack}{RGB}{50,50,50}

% --- Difficulty badges (rounded corners via tikz, pre-rendered in saveboxes) ---
\usepackage{tikz}
\newsavebox{\badgechubox}
\savebox{\badgechubox}{\tikz[baseline]{\node[fill=badgegreen,text=white,rounded corners=2.5pt,inner xsep=4pt,inner ysep=1.5pt,font=\sffamily\footnotesize\bfseries]{初级};}}
\newsavebox{\badgezhongbox}
\savebox{\badgezhongbox}{\tikz[baseline]{\node[fill=badgeblue,text=white,rounded corners=2.5pt,inner xsep=4pt,inner ysep=1.5pt,font=\sffamily\footnotesize\bfseries]{中级};}}
\newsavebox{\badgegaobox}
\savebox{\badgegaobox}{\tikz[baseline]{\node[fill=badgeblack,text=white,rounded corners=2.5pt,inner xsep=4pt,inner ysep=1.5pt,font=\sffamily\footnotesize\bfseries]{高级};}}
\newcommand{\lvchu}{\usebox{\badgechubox}}
\newcommand{\lvzhong}{\usebox{\badgezhongbox}}
\newcommand{\lvgao}{\usebox{\badgegaobox}}
\pdfstringdefDisableCommands{%
  \def\lvchu{[初级]}%
  \def\lvzhong{[中级]}%
  \def\lvgao{[高级]}%
}

\lstset{
  basicstyle=\ttfamily\scriptsize,
  keywordstyle=\color{codekw}\bfseries,
  commentstyle=\color{codecomment}\itshape,
  stringstyle=\color{codestring},
  backgroundcolor=\color{codebg},
  frame=single,
  rulecolor=\color{codeframe},
  breaklines=true,
  breakatwhitespace=false,
  breakindent=0pt,
  postbreak=\mbox{\textcolor{red}{$\hookrightarrow$}\space},
  numbers=left,
  numberstyle=\tiny\color{gray},
  columns=flexible,
  keepspaces=true,
  showstringspaces=false,
  tabsize=2,
  aboveskip=4pt,
  belowskip=4pt,
  extendedchars=true,
  literate={，}{{，}}1 {；}{{；}}1 {！}{{！}}1 {？}{{？}}1 {“}{{“}}1 {”}{{”}}1 {（}{{（}}1 {）}{{）}}1 {《}{{《}}1 {》}{{》}}1,
}

% --- Extra language definitions (no built-in listings support) ---
\lstdefinelanguage{json}{
  morekeywords={true,false,null},
  sensitive=true,
  morestring=[b]",
  comment=[l]{//},
  morecomment=[s]{/*}{*/},
}
\lstdefinelanguage{yaml}{
  morekeywords={true,false,null,yes,no,on,off},
  sensitive=false,
  morecomment=[l]{\#},
  morestring=[b]",
  morestring=[d]',
}
\lstdefinelanguage{http}{
  morekeywords={GET,POST,PUT,DELETE,PATCH,HEAD,OPTIONS,HTTP,Host,Accept,Authorization,Content-Type,Content-Length,User-Agent},
  sensitive=true,
  morecomment=[l]{\#},
  morestring=[b]",
}
\lstdefinelanguage{TypeScript}{
  morekeywords={abstract,any,as,async,await,break,case,catch,class,const,continue,debugger,declare,default,delete,do,else,enum,export,extends,false,finally,for,from,function,get,if,implements,import,in,instanceof,interface,keyof,let,module,namespace,never,new,null,number,of,package,private,protected,public,readonly,return,set,static,string,super,switch,symbol,this,throw,true,try,type,typeof,undefined,unknown,var,void,while,with,yield,Record,Array,Promise,AsyncIterable,ReadableStream,Date,Error,Object,JSON},
  sensitive=true,
  morecomment=[l]{//},
  morecomment=[s]{/*}{*/},
  morestring=[b]",
  morestring=[b]',
  morestring=[b]`{`},
}

% --- Compact spacing ---
\linespread{0.95}
\setlength{\parskip}{1pt plus 1pt minus 1pt}
\setlength{\parindent}{0pt}

% --- Table styling ---
\renewcommand{\arraystretch}{1.1}

% --- Heading styling (numbered) ---
\titleformat{\section}{\normalsize\bfseries}{\thesection\quad}{0em}{}
\titlespacing{\section}{0pt}{6pt}{2pt}
\titleformat{\subsection}{\small\bfseries}{\thesubsection\quad}{0em}{}
\titlespacing{\subsection}{0pt}{4pt}{1pt}
\titleformat{\subsubsection}{\small\bfseries}{\thesubsubsection\quad}{0em}{}
\titlespacing{\subsubsection}{0pt}{3pt}{1pt}

% --- Page style ---
\pagestyle{fancy}
\fancyhf{}
\fancyhead[L]{\small\emph{\leftmark}}
\fancyhead[R]{\small\thepage}
\renewcommand{\headrulewidth}{0.4pt}
\renewcommand{\sectionmark}[1]{}  % prevent sections from overwriting leftmark

% --- Hyperref setup ---
\hypersetup{
  colorlinks=true,
  linkcolor=blue,
  urlcolor=blue,
  citecolor=blue,
}

% --- Caption format ---
\DeclareCaptionFormat{myformat}{\small\bfseries #1#2#3}
\captionsetup{format=myformat}

\begin{document}

\title{大模型基础面试题总结}
\date{}
\maketitle
\markboth{Java 版 Agent 知识}{ Java 版 Agent 知识}

很多同学准备大模型面试时，第一反应是去背 Transformer、Attention、RLHF 这些词。不是说这些不重要，但对大部分后端转 AI 应用开发、AI 工程应用岗位来说，面试官更关心的是另一件事：



\textbf{你是不是真的理解大模型调用链路里的工程约束。}



比如 Token 为什么会影响成本和延迟？上下文窗口为什么不是越大越好？Temperature 为什么会影响结构化输出稳定性？Function Calling 为什么不能让模型直接执行真实业务操作？这些问题看起来基础，答不好就会暴露一个信号：你可能只是调过 API，还没有把大模型当作生产系统里的一个不稳定外部依赖来治理。



这份大模型基础面试题主要根据 AI 专栏现有文章整理。它不是让你机械背题，而是帮你建立一条复习主线：


\begin{enumerate}[itemsep=1pt, topsep=2pt]
  \item 先理解 \textbf{Token、上下文窗口、采样参数}，知道模型为什么会不稳定。
  \item 再理解 \textbf{API 调用工程}，知道一次模型调用在生产里要经过哪些治理环节。
  \item 接着理解 \textbf{结构化输出与工具调用}，知道怎么让模型输出能被程序消费。
  \item 最后理解 \textbf{AI 应用评测}，知道怎么判断你的 AI 应用到底有没有变好。
\end{enumerate}


\section*{面试官真正想考什么}


大模型基础题表面上问概念，实际考的是工程判断。你可以按下面这张表来理解。


{\footnotesize
\begin{tabular}{|p{\dimexpr 0.159\columnwidth-2\tabcolsep\relax}|p{\dimexpr 0.365\columnwidth-2\tabcolsep\relax}|p{\dimexpr 0.476\columnwidth-2\tabcolsep\relax}|}
\hline
\textbf{考察方向} & \textbf{面试官想确认什么} & \textbf{常见扣分点} \\
\hline
Token 和上下文 & 你是否理解成本、延迟、窗口限制和信息取舍 & 只说 Token 是“词元”，讲不出工程影响 \\
\hline
采样参数 & 你是否知道如何在创造性和稳定性之间取舍 & 把 Temperature 说成越高越聪明 \\
\hline
API 调用链路 & 你是否具备把模型接入生产系统的经验 & 只说调用 HTTP 接口，忽略重试、限流、幂等 \\
\hline
结构化输出 & 你是否知道自然语言约束不等于工程契约 & 认为“请返回 JSON”就足够可靠 \\
\hline
评测闭环 & 你是否能验证效果，而不是凭感觉调 Prompt & 只看公开 benchmark，不做业务 Golden Set \\
\hline
\end{tabular}
}



一个不错的回答通常不是定义式的，而是“概念 + 问题 + 工程解法”。例如问 Token，你可以先解释 Token 是模型处理文本的基本单位，再补一句：Token 直接影响上下文容量、推理成本、响应延迟和截断风险，所以生产系统里要做预算估算、历史消息压缩、RAG 证据筛选和最大输出限制。



这就比单纯背定义强很多。



\section*{LLM 运行机制}


参考文章：《LLM 运行机制：Token、上下文窗口与采样参数怎么影响输出》



这一组题是大模型面试的地基。不要只记术语，要重点理解这些概念如何影响真实系统的稳定性、成本和答案质量。



建议掌握这些关键点：


\begin{itemize}[itemsep=1pt, topsep=2pt]
  \item Token 不是字符，也不是中文里的“字”。不同语言、符号、代码片段的切分方式不同，因此同样长度的中文、英文、代码，Token 消耗可能差很多。
  \item 上下文窗口不是无限记忆。窗口越大，成本、延迟、噪声、Lost in the Middle 风险都会增加。
  \item Temperature、Top-P、Top-K 控制的是采样分布，不是模型“智商”。生产环境通常更关注稳定性和可复现性。
  \item 幻觉不是单靠某个参数就能消灭的。更可靠的做法是 RAG、工具调用、引用来源、输出校验和评测闭环一起做。
\end{itemize}


高频面试题：


\begin{itemize}[itemsep=1pt, topsep=2pt]
  \item Token 是什么？为什么中文、英文、代码消耗的 Token 不一样？
  \item 上下文窗口是什么？上下文窗口越大，效果一定越好吗？
  \item 什么是 Lost in the Middle 问题？长上下文场景下怎么缓解？
  \item Temperature、Top-P、Top-K 分别控制什么？生产环境怎么设置更稳？
  \item 为什么 Temperature 设置为 0，模型输出仍然可能不完全一致？
  \item 大模型为什么会产生幻觉？常见缓解方案有哪些？
  \item Token 预算怎么估算？输入、输出、历史消息、RAG 证据如何取舍？
  \item 长上下文窗口会不会取代 RAG？二者分别适合什么场景？
\end{itemize}


面试追问通常会落到场景上。比如“你们的客服机器人历史会话太长怎么办？”这时不要只说“做摘要”，更完整的回答是：先区分必须保留的业务状态、最近对话、用户画像和可丢弃闲聊；再做 Token 预算；超过阈值时对历史消息做结构化摘要；RAG 证据只放最相关片段；最后通过评测集验证压缩后是否影响关键问题回答。



\begin{figure}[H]
\centering
\includegraphics[width=0.85\textwidth]{/home/zhang/demos/latex_demos/Agent知识/HTML/大模型基础面试题总结 _ JavaGuide_files/llm-token-process.jpg}
\caption*{Token 化过程示例}
\end{figure}



\section*{API 调用工程}


参考文章：《大模型 API 调用工程实践：流式输出、重试、限流与结构化返回》



这一组题考的是你有没有把模型当作生产依赖来治理。大模型 API 和普通 HTTP API 很像，但又更麻烦：它慢、贵、不稳定、输出不可完全控，还可能被供应商限流。



建议掌握这些关键点：


\begin{itemize}[itemsep=1pt, topsep=2pt]
  \item 一次模型调用不只是“发请求拿结果”，而是一条完整链路：请求校验、Prompt 组装、上下文注入、模型路由、限流、超时、重试、流式返回、结构化解析、日志和评测。
  \item Streaming 主要改善首字体验，不等于减少总耗时，也不等于降低 Token 成本。
  \item 重试必须和幂等绑定。没有幂等设计，重试可能造成重复扣费、重复落库、重复执行工具。
  \item 限流不能只看 QPS，还要看 RPM、TPM、并发数、上下文大小、最大输出和租户预算。
\end{itemize}


高频面试题：


\begin{itemize}[itemsep=1pt, topsep=2pt]
  \item 大模型 API 调用的完整链路是什么？
  \item Streaming 为什么能改善用户体验？它能减少总耗时和 Token 成本吗？
  \item SSE、WebSocket、HTTP Chunked 在流式输出场景下怎么选？
  \item 哪些大模型 API 错误可以重试？哪些错误不能重试？
  \item 为什么大模型调用必须做幂等？
  \item 大模型限流为什么不能只按 QPS 做？
  \item 模型网关通常要承担哪些能力？
  \item AI 应用的调用日志里至少要记录哪些字段？
\end{itemize}


一个比较稳的回答方式是先讲“链路”，再讲“治理”。例如回答“为什么需要模型网关”，可以这样展开：模型网关把供应商差异、模型路由、fallback、限流、熔断、Token 预算、成本归因和观测统一起来，避免业务代码直接耦合某个模型供应商。业务只关心能力，网关负责稳定性和成本。



\section*{结构化输出与工具调用}


参考文章：《大模型结构化输出：从 JSON 契约到 Function Calling 落地》



这一组题是 AI 应用开发的高频追问点。因为只要模型输出要进业务系统，就绕不开结构化输出、Schema 校验和工具调用安全。



建议掌握这些关键点：


\begin{itemize}[itemsep=1pt, topsep=2pt]
  \item “请返回 JSON”只是自然语言提示，不是强约束。模型可能多输出解释、漏字段、类型错误、枚举乱写。
  \item JSON Mode 主要保证合法 JSON，Structured Outputs 更关注是否符合 Schema，但服务端仍然必须校验。
  \item Function Calling 的本质是让模型生成工具调用意图，真正执行权在业务系统。
  \item MCP 解决的是工具如何标准化接入宿主，Function Calling 解决的是模型如何表达调用意图，它们不在同一层。
  \item 工具调用必须做参数校验、权限校验、二次确认、幂等、审计和超时控制。
\end{itemize}


高频面试题：


\begin{itemize}[itemsep=1pt, topsep=2pt]
  \item 为什么只写“请返回 JSON”不可靠？
  \item JSON Mode 和 Structured Outputs 有什么区别？
  \item JSON Schema 在大模型应用里解决什么问题？
  \item Function Calling 的完整链路是什么？
  \item Function Calling 和 MCP 有什么区别？
  \item MCP Tool 和普通 HTTP API 有什么关系？
  \item Agent Skill 和 Function Calling 是一回事吗？
  \item 结构化输出失败后怎么处理？
  \item 工具调用为什么必须做安全治理？
  \item 面试里怎么一句话概括结构化输出？
\end{itemize}


这类题最容易答得太抽象。建议始终带一个业务例子：比如“退款工具调用”。模型可以生成 \texttt{refundOrder(orderId, amount, reason)} 的调用参数，但后端必须确认当前用户是否有权限、订单是否属于本人、金额是否可退、是否已经退过、是否需要二次确认。模型只能提出意图，不能绕过业务规则。



\section*{AI 应用评测}


参考文章：《AI 应用评测体系：从 Golden Set 构建到线上灰度闭环》



很多候选人会调 Prompt，但说不清“怎么证明调得更好了”。这就是评测题的价值。面试官问评测，通常是在判断你有没有生产意识。



建议掌握这些关键点：


\begin{itemize}[itemsep=1pt, topsep=2pt]
  \item 公开 benchmark 只能粗略判断模型通用能力，不能代表你的业务数据分布。
  \item Golden Set 的价值不在数量，而在分布。正常路径、边缘场景、对抗样本、高权重失败都要覆盖。
  \item LLM-as-Judge 可以提高评测效率，但有位置偏差、冗长偏差、同源偏差和推理能力边界，不能完全替代人工。
  \item RAG 和 Agent 都要分段评测。只看最终答案，很难定位问题来自检索、生成、工具调用还是执行轨迹。
\end{itemize}


高频面试题：


\begin{itemize}[itemsep=1pt, topsep=2pt]
  \item 为什么不能只靠公开 benchmark 评估 AI 应用质量？
  \item Golden Set 应该怎么构建？冷启动阶段没有生产日志怎么办？
  \item LLM-as-Judge 有哪些主要偏差？怎么缓解？
  \item RAG 评测为什么必须分检索和生成两段？
  \item Agent 评测为什么比普通问答和 RAG 更复杂？
  \item 离线评测、Trace 回放、线上灰度分别解决什么问题？
  \item CI 里的 AI 评测如何平衡速度和覆盖度？
  \item 如果 LLM-as-Judge 和人工评测结果不一致，应该怎么处理？
\end{itemize}


回答评测题时，尽量形成闭环：先有 Golden Set 做离线回归，再用 Trace 回放覆盖真实线上路径，最后通过灰度和线上采样验证真实用户分布。没有这条链路，优化基本靠感觉。



\section*{答题框架}


大模型基础题可以套用一个简单框架：


\begin{enumerate}[itemsep=1pt, topsep=2pt]
  \item 先解释概念：用一句话说清楚它是什么。
  \item 再说明影响：它会影响质量、成本、延迟、稳定性还是安全。
  \item 接着给工程做法：生产里如何配置、校验、降级或观测。
  \item 最后补充边界：在哪些场景下会失效，或者需要和其他方案组合。
\end{enumerate}


比如问“长上下文会不会取代 RAG”，可以这样答：



长上下文能提升单次输入容量，适合少量文档的深度分析，但它不能完全取代 RAG。企业知识库通常有海量文档、权限隔离、频繁更新、成本控制和引用溯源要求，不可能每次把所有内容塞进窗口。更现实的做法是用 RAG 做候选证据筛选，再把少量高质量上下文交给长上下文模型处理。



\section*{常见扣分点}

\begin{itemize}[itemsep=1pt, topsep=2pt]
  \item 只背定义，不讲工程影响。
  \item 把大模型 API 当普通 HTTP 接口，没有限流、重试、幂等、观测意识。
  \item 认为结构化输出等于“让模型返回 JSON”，忽略 Schema 和服务端校验。
  \item 认为 Function Calling 是模型直接执行函数，忽略业务系统的执行权和安全边界。
  \item 只看模型排行榜，不知道 Golden Set、Trace 回放和线上灰度。
\end{itemize}


\section*{复习建议}


如果时间有限，建议按这个顺序复习：


\begin{enumerate}[itemsep=1pt, topsep=2pt]
  \item 先看 Token、上下文窗口、采样参数，建立基础认知。
  \item 再看 API 调用工程，理解从 Demo 到生产的差距。
  \item 接着看结构化输出和 Function Calling，这是 AI 应用开发的高频追问点。
  \item 最后看评测体系，尤其是 Golden Set、LLM-as-Judge、Trace 回放。
\end{enumerate}


复习时不要只问自己“这个概念是什么”，还要继续追问三句：生产里会出什么问题？怎么定位？怎么治理？能答到这个层次，大模型基础面试基本就稳了。


\end{document}
